% Relativistic Chapter XII, §110: Axial Magnetic Field on Gravitational Cylinder Instability
% Relativistic generalization of Chandrasekhar, Chapter XII §110
% Conventions: see RELATIVISTIC_CONVENTIONS.md

\chapter*{Relativistic Gravitational Instability of a Cylinder with an Axial Magnetic Field}
\addcontentsline{toc}{chapter}{Relativistic Gravitational Instability of a Cylinder with Axial $B$-Field}
\label{ch:rel-12-110}

%% ================================================================
\section{Introduction: magnetized relativistic self-gravitating cylinders}
\label{sec:rel-12-110-intro}

In the non-relativistic treatment of \S~110, Chandrasekhar showed that
a uniform axial magnetic field partially stabilizes an infinite
self-gravitating cylinder: it raises the minimum wavelength at which
the gravitational (Jeans-type) varicose instability can operate, but
can never suppress the instability entirely---modes of sufficiently
long wavelength always remain unstable, owing to the logarithmic
divergence of the gravitational potential perturbation.

The present section extends that analysis to the relativistic regime,
where the inertia of the fluid is governed by the relativistic
enthalpy density $w = (\varepsilon + p)/c^{2}$ rather than the
rest-mass density $\rho_{0}$, and the magnetic field itself
contributes to the stress--energy tensor.  This generalization is
relevant in several astrophysical contexts:

\begin{enumerate}
\item \textbf{Relativistic jet fragmentation.}  Extragalactic jets,
  launched with bulk Lorentz factors $\Gamma \sim 10$--$50$ from
  active galactic nuclei, carry helical magnetic fields whose axial
  component may suppress varicose (sausage-mode) kink instabilities
  that would otherwise fragment the jet.  Whether such a jet fragments
  into knots depends on the competition between relativistic
  self-gravity and magnetic tension---precisely the physics of the
  relativistic \S~110 problem.
\item \textbf{Magnetized neutron-star filaments.}  Dense magnetized
  filaments in neutron-star magnetospheres or in the post-merger
  ejecta of binary neutron-star coalescences can be relativistically
  hot, with $p/\varepsilon$ not negligible.
\item \textbf{Cosmic string analogues.}  Gravitationally bound
  magnetized plasma columns in the early universe or in gamma-ray
  burst engines involve relativistic equation-of-state effects.
\end{enumerate}

Throughout, we work in flat Minkowski spacetime with cylindrical
coordinates $(t, \varpi, \varphi, z)$ and metric signature
$(-,+,+,+)$.  The speed of light $c$ is kept explicit.  The
four-velocity satisfies $u^{\mu}u_{\mu} = -c^{2}$.  We employ Gaussian
units for electromagnetic quantities, consistent with the original
treatment.

%% ================================================================
\section{Relativistic MHD framework for an axially magnetized cylinder}
\label{sec:rel-12-110-framework}

\subsection{Equilibrium configuration}
\label{sec:rel-12-110-equil}

Consider an infinite cylinder of radius $R$, filled with a
relativistic perfect fluid of energy density $\varepsilon$, pressure
$p$, and rest-mass density $\rho_{0}$.  A uniform magnetic field
$\mathbf{H} = H\,\hat{e}_{z}$ threads the cylinder along its axis.
In the fluid rest frame the magnetic field four-vector is
\begin{equation}
\label{eq:rel-12-110-1}
b^{\mu} = \bigl(0,\,0,\,0,\,H\bigr),\qquad
b^{2} \equiv b^{\mu}b_{\mu} = H^{2},
\end{equation}
since the fluid is at rest in the equilibrium state.

The total energy--momentum tensor for the magnetized fluid is
(cf.\ RELATIVISTIC\_CONVENTIONS.md)
\begin{equation}
\label{eq:rel-12-110-2}
T^{\mu\nu} = \Bigl(\varepsilon + p + \frac{b^{2}}{4\pi}\Bigr)
\frac{u^{\mu}u^{\nu}}{c^{2}}
+ \Bigl(p + \frac{b^{2}}{8\pi}\Bigr) g^{\mu\nu}
- \frac{b^{\mu}b^{\nu}}{4\pi}.
\end{equation}
The effective relativistic enthalpy density (including the magnetic
contribution) is
\begin{equation}
\label{eq:rel-12-110-3}
w_{\mathrm{tot}} = \frac{1}{c^{2}}\Bigl(\varepsilon + p
+ \frac{b^{2}}{4\pi}\Bigr)
= \frac{1}{c^{2}}\Bigl(\varepsilon + p
+ \frac{H^{2}}{4\pi}\Bigr).
\end{equation}
In the non-relativistic limit $\varepsilon \to \rho_{0}c^{2}$,
$p \ll \rho_{0}c^{2}$, $H^{2}/(8\pi) \ll \rho_{0}c^{2}$, we recover
$w_{\mathrm{tot}} \to \rho_{0}$.

\subsection{Relativistic Alfv\'en speed}
\label{sec:rel-12-110-alfven}

A central quantity in the relativistic generalization is the
Alfv\'en speed.  For a magnetic field aligned along $\hat{e}_{z}$
the relativistic Alfv\'en speed is (cf.\ RELATIVISTIC\_CONVENTIONS.md)
\begin{equation}
\label{eq:rel-12-110-4}
\boxed{\vA^{2}
= \frac{H^{2}/(4\pi)}{(\varepsilon + p)/c^{2} + H^{2}/(4\pi c^{2})}
= \frac{H^{2} c^{2}}{4\pi(\varepsilon + p) + H^{2}}.}
\end{equation}
\paragraph{Causality.}
One verifies immediately that $\vA^{2} < c^{2}$ for all finite values
of $H$, $\varepsilon$, and $p > 0$:
\begin{equation}
\label{eq:rel-12-110-5}
\vA^{2} = \frac{c^{2}}{1 + 4\pi(\varepsilon + p)/H^{2}} < c^{2}.
\end{equation}
The non-relativistic limit gives
$\vA^{2} \to \mu H^{2}/(4\pi\rho_{0})$ (with $\mu=1$ in Gaussian
units), recovering the classical Alfv\'en frequency
$\Omega_{a}^{2} = k^{2}\vA^{2}$ of equation~(115) in \S~110.

%% ================================================================
\section{Perturbation equations: relativistic axisymmetric modes}
\label{sec:rel-12-110-perturb}

We restrict attention to axisymmetric perturbations, as in the
non-relativistic case only these modes are gravitationally unstable.
Let the boundary of the cylinder be deformed as
\begin{equation}
\label{eq:rel-12-110-6}
\varpi = R + \epsilon_{0}\,e^{ikz + \sigma t},
\end{equation}
where $k$ is the axial wave number and $\sigma$ the (possibly complex)
growth rate.  We define $x = kR$.

In the ideal MHD limit (infinite conductivity, $F^{\mu\nu}u_{\nu}=0$
in the fluid frame), the linearized relativistic momentum equation
for the perturbed four-velocity $\delta u^{\mu}$ in the rest frame of
the equilibrium fluid becomes
\begin{equation}
\label{eq:rel-12-110-7}
w_{\mathrm{tot}}\,\sigma\,\delta u^{i}
= -\partial^{i}\delta\Pi_{\mathrm{rel}}
+ \frac{H}{4\pi}\,\frac{\partial\,\delta b^{i}}{\partial z},
\end{equation}
where $\delta\Pi_{\mathrm{rel}}$ is the relativistic generalization of
the total (gravitational + fluid + magnetic) pressure perturbation:
\begin{equation}
\label{eq:rel-12-110-8}
\delta\Pi_{\mathrm{rel}} = -\delta\mathfrak{B}
+ \frac{\delta p}{w_{\mathrm{tot}}}
+ \frac{H\,\delta b_{z}}{4\pi\,w_{\mathrm{tot}}}.
\end{equation}
Here $\delta\mathfrak{B}$ denotes the perturbation in the
gravitational potential.  The induction equation retains its classical
form in the fluid frame:
\begin{equation}
\label{eq:rel-12-110-9}
\sigma\,\delta b^{i} = ik\,H\,\delta u^{i}.
\end{equation}

Equation~\eqref{eq:rel-12-110-7} can be written in terms of a scalar
function $U(\varpi)$ exactly as in the non-relativistic case (since
the fluid is at rest in equilibrium and the perturbations are
non-relativistic in the linear regime).  Substituting
\eqref{eq:rel-12-110-9} into \eqref{eq:rel-12-110-7}, one obtains the
eigenvalue equation
\begin{equation}
\label{eq:rel-12-110-10}
\Bigl(\sigma^{2} + \Omega_{a,\mathrm{rel}}^{2}\Bigr) U
= -\frac{\epsilon_{0}\Pi_{0}}{w_{\mathrm{tot}}}\,k\,I_{0}(k\varpi),
\end{equation}
where we have defined the \emph{relativistic Alfv\'en frequency}
\begin{equation}
\label{eq:rel-12-110-11}
\boxed{\Omega_{a,\mathrm{rel}}^{2}
= k^{2}\vA^{2}
= \frac{k^{2} H^{2} c^{2}}{4\pi(\varepsilon+p) + H^{2}}.}
\end{equation}
In the non-relativistic limit,
$\Omega_{a,\mathrm{rel}}^{2} \to k^{2}\mu H^{2}/(4\pi\rho_{0})
= \Omega_{a}^{2}$, recovering equation~(115) of \S~110.

%% ================================================================
\section{Dispersion relation: relativistic characteristic equation}
\label{sec:rel-12-110-dispersion}

The boundary conditions---kinematic compatibility at $\varpi = R$,
continuity of the normal component of $\mathbf{b}$ across $\varpi=R$,
and continuity of the total normal stress on the deformed
surface---follow the same algebraic steps as in \S~110, with the
replacement $\rho \to w_{\mathrm{tot}} = (\varepsilon + p + H^{2}/4\pi)/c^{2}$.
The gravitational potential perturbation, however, is sourced by the
\emph{total} energy density $\varepsilon$ (not merely $\rho_{0}$) in
general relativity.  For a weak-field (linearized GR) treatment valid
to post-Newtonian order, the Poisson equation is replaced by
\begin{equation}
\label{eq:rel-12-110-12}
\nabla^{2}\delta\mathfrak{B}
= -4\pi G\,\frac{1}{c^{2}}\bigl(\delta\varepsilon
+ 3\,\delta p + \delta b^{2}/4\pi\bigr),
\end{equation}
but for the dominant gravitational mode on scales much larger than the
Schwarzschild radius of the cylinder, we may approximate the source
term by $-4\pi G\,\rho_{\mathrm{eff}}$, where the effective
gravitational mass density is
\begin{equation}
\label{eq:rel-12-110-13}
\rho_{\mathrm{eff}} = \frac{\varepsilon + p}{c^{2}}
+ \frac{H^{2}}{8\pi c^{2}} + \frac{3p}{c^{2}}
\approx \frac{\varepsilon + p}{c^{2}}\Bigl(1 + \frac{2p}{\varepsilon+p}\Bigr)
+ \frac{H^{2}}{8\pi c^{2}}.
\end{equation}
For simplicity and transparency of the relativistic corrections, we
define the \emph{relativistic gravitational mass density}
\begin{equation}
\label{eq:rel-12-110-14}
\rho_{G} \equiv \frac{\varepsilon + 3p}{c^{2}}
+ \frac{H^{2}}{8\pi c^{2}},
\end{equation}
which reduces to $\rho_{0}$ in the Newtonian limit.

Proceeding exactly as in the non-relativistic derivation---applying
boundary conditions (1)--(3) of \S~110 with the replacements
$\rho \to w_{\mathrm{tot}}$ (inertial) and $\rho \to \rho_{G}$
(gravitational source)---we arrive at the \textbf{relativistic
characteristic equation}:
\begin{equation}
\label{eq:rel-12-110-15}
\boxed{\sigma^{2}
= 4\pi G\,\rho_{G}\,\frac{x\,I_{1}(x)}{I_{0}(x)}
\Bigl[K_{0}(x)\,I_{0}(x) - \tfrac{1}{2}\Bigr]
- \Omega_{a,\mathrm{rel}}^{2}\,
\frac{x^{2}\,I_{1}(x)\,K_{0}(x)}{I_{0}(x)\,K_{1}(x)},}
\end{equation}
where $x = kR$ and $I_{n}$, $K_{n}$ are modified Bessel functions.
Defining the relativistic critical field strength
\begin{equation}
\label{eq:rel-12-110-16}
H_{0,\mathrm{rel}} \equiv 4\pi R\,c\,
\sqrt{\frac{\rho_{G}}{w_{\mathrm{tot}}}}\,\sqrt{G\,\rho_{G}},
\end{equation}
we can recast the dispersion relation~\eqref{eq:rel-12-110-15} in the
dimensionless form
\begin{equation}
\label{eq:rel-12-110-17}
\frac{\sigma^{2}}{4\pi G\,\rho_{G}}
= \frac{x\,I_{1}(x)}{I_{0}(x)}
\left\{K_{0}(x)\,I_{0}(x) - \tfrac{1}{2}
- \Bigl(\frac{H}{H_{0,\mathrm{rel}}}\Bigr)^{2}
\frac{x^{2}\,K_{0}(x)}{I_{0}(x)\,K_{1}(x)}\right\}.
\end{equation}
This is the direct relativistic generalization of equation~(128) of
\S~110.

\paragraph{Non-relativistic limit.}
Setting $\varepsilon = \rho_{0}c^{2}$, $p\to 0$,
$H^{2}/(8\pi) \ll \rho_{0}c^{2}$, we find
$\rho_{G} \to \rho_{0}$, $w_{\mathrm{tot}} \to \rho_{0}$,
$\Omega_{a,\mathrm{rel}}^{2} \to \mu H^{2}k^{2}/(4\pi\rho_{0})$,
$H_{0,\mathrm{rel}} \to 4\pi\rho_{0}R\sqrt{G\rho_{0}} = H_{0}$, and
equation~\eqref{eq:rel-12-110-17} reduces exactly to the classical
equation~(128).

%% ================================================================
\section{Nature of the stabilizing effect: relativistic magnetic tension}
\label{sec:rel-12-110-nature}

\subsection{Minimum wavelength for instability}
\label{sec:rel-12-110-minstab}

The structure of equation~\eqref{eq:rel-12-110-17} shows that the
axial magnetic field contributes a stabilizing term proportional to
$H^{2}$ (through $\Omega_{a,\mathrm{rel}}^{2}$), which \emph{grows}
with wave number $k$.  Meanwhile, the gravitational term has the
long-wavelength ($x\to 0$) logarithmic divergence
$K_{0}(x)\,I_{0}(x) - \frac{1}{2} \to -\log(x/2) - \gamma_{E}$,
where $\gamma_{E}$ is the Euler--Mascheroni constant.  Consequently,
as in the non-relativistic case:

\begin{quote}
\emph{No axial magnetic field, however strong, can stabilize the
cylinder against perturbations of all wavelengths.  Modes of
sufficiently long wavelength remain gravitationally unstable.}
\end{quote}

This fundamental conclusion is unchanged by relativistic corrections,
because the logarithmic divergence of the gravitational potential is a
purely geometric (Newtonian-order) effect that persists in linearized
GR.

The critical wave number $x_{a}$ below which instability occurs is the
root of
\begin{equation}
\label{eq:rel-12-110-18}
K_{0}(x)\,I_{0}(x) - \tfrac{1}{2}
= \Bigl(\frac{H}{H_{0,\mathrm{rel}}}\Bigr)^{2}
\frac{x^{2}\,K_{0}(x)}{I_{0}(x)\,K_{1}(x)}.
\end{equation}
For $\sigma^{2} < 0$ when $x > x_{a}$ and $\sigma^{2} > 0$ for
$0 < x < x_{a}$, the cylinder is unstable to all varicose
deformations with wavelengths exceeding
\begin{equation}
\label{eq:rel-12-110-19}
\lambda_{\min} = \frac{2\pi R}{x_{a}}.
\end{equation}

\subsection{Critical field strength: $B^{2}/(8\pi)$ vs
$(\varepsilon+p)\times$gravitational energy}
\label{sec:rel-12-110-critical}

The competition between magnetic tension and gravitational collapse is
governed by the ratio $H/H_{0,\mathrm{rel}}$.  The critical field
$H_{0,\mathrm{rel}}$ can be written as
\begin{equation}
\label{eq:rel-12-110-20}
\frac{H_{0,\mathrm{rel}}^{2}}{8\pi}
= 2\pi\,R^{2}\,G\,\rho_{G}^{2}\,\frac{c^{2}}{w_{\mathrm{tot}}c^{2}}
= 2\pi\,R^{2}\,G\,\frac{\rho_{G}^{2}}{w_{\mathrm{tot}}}.
\end{equation}
In the non-relativistic limit this gives
$H_{0}^{2}/(8\pi) = 2\pi R^{2} G \rho_{0}^{2}$, which is the
magnetic pressure that equals the gravitational binding energy per
unit volume of the cylinder.

In the relativistic regime, the stabilization criterion becomes
\begin{equation}
\label{eq:rel-12-110-21}
\frac{H^{2}}{8\pi} \gtrsim 2\pi R^{2}\,G\,
\frac{(\varepsilon + 3p)^{2}/c^{4}}{(\varepsilon + p + H^{2}/4\pi)/c^{2}}
\quad\Longrightarrow\quad
\text{substantial suppression of short-wavelength modes.}
\end{equation}
The key relativistic effects are:
\begin{enumerate}
\item \textbf{Enhanced gravitational source:} the factor
  $\rho_{G} = (\varepsilon + 3p + H^{2}/8\pi)/c^{2}$ replaces
  $\rho_{0}$, so relativistic pressure (and the magnetic field itself)
  contribute to the gravitational instability drive---an effect absent
  in Newtonian theory.
\item \textbf{Enhanced inertia:}  the effective inertial mass
  $w_{\mathrm{tot}}$ includes the enthalpy $(\varepsilon+p)$ and the
  magnetic energy density $H^{2}/(4\pi)$, which slows the growth of
  perturbations.
\item \textbf{Bounded Alfv\'en speed:}  $\vA < c$ always, preventing
  any superluminal signal propagation in the stabilization mechanism.
\end{enumerate}

\subsection{Strong-field asymptotic behaviour}
\label{sec:rel-12-110-strongB}

For $H \gg H_{0,\mathrm{rel}}$, the critical wavenumber $x_{a} \ll 1$
and we may use the small-argument expansions of the Bessel functions.
The analysis parallels equations~(132)--(138) of \S~110, and we obtain
\begin{equation}
\label{eq:rel-12-110-22}
x_{a} \simeq 2\exp\!\Bigl[-\gamma_{E} - \tfrac{1}{2}
- 2\Bigl(\frac{H}{H_{0,\mathrm{rel}}}\Bigr)^{2}\Bigr]
= 0.6811\,\exp\!\Bigl[-2\Bigl(\frac{H}{H_{0,\mathrm{rel}}}\Bigr)^{2}\Bigr],
\end{equation}
\begin{equation}
\label{eq:rel-12-110-23}
x_{m} \simeq 0.4134\,\exp\!\Bigl[-2\Bigl(\frac{H}{H_{0,\mathrm{rel}}}\Bigr)^{2}\Bigr],
\end{equation}
and the maximum growth rate
\begin{equation}
\label{eq:rel-12-110-24}
\sigma_{\max}^{2}
\simeq \frac{4\pi G\,\rho_{G}}{4e^{2}}\,x_{m}^{2}
= 0.2055\,(4\pi G\,\rho_{G})\,
\exp\!\Bigl[-4\Bigl(\frac{H}{H_{0,\mathrm{rel}}}\Bigr)^{2}\Bigr].
\end{equation}
The exponential suppression of the growth rate with increasing
$H/H_{0,\mathrm{rel}}$ is identical in form to the non-relativistic
result, but the \emph{scale} is set by $\rho_{G}$ rather than
$\rho_{0}$, leading to faster growth in relativistic configurations
(since $\rho_{G} > \rho_{0}$ when $p > 0$).

%% ================================================================
\section{Causality constraint: $\vA < c$}
\label{sec:rel-12-110-causality}

\subsection{The Alfv\'en speed bound}
\label{sec:rel-12-110-vabound}

From equation~\eqref{eq:rel-12-110-4}, the relativistic Alfv\'en
speed satisfies
\begin{equation}
\label{eq:rel-12-110-25}
\frac{\vA^{2}}{c^{2}}
= \frac{H^{2}}{4\pi(\varepsilon + p) + H^{2}} < 1
\end{equation}
for all $H < \infty$, provided $\varepsilon + p > 0$ (satisfied by
any physical equation of state).  In the extreme magnetic-dominance
limit $H^{2}/(4\pi) \gg \varepsilon + p$, one has $\vA \to c$,
approaching but never reaching the speed of light.  This ensures that
the stabilization mechanism cannot transmit information superluminally.

\subsection{Fast magnetosonic speed}
\label{sec:rel-12-110-fastms}

The fast magnetosonic speed, which governs compressive perturbations
coupled to the magnetic field, is
\begin{equation}
\label{eq:rel-12-110-26}
v_{f}^{2} = \cs^{2} + \vA^{2}
- \frac{\cs^{2}\,\vA^{2}}{c^{2}},
\end{equation}
where $\cs^{2} = (\partial p/\partial\varepsilon)_{s}$ is the sound
speed.  One verifies that $v_{f}^{2} < c^{2}$ provided
$\cs^{2} < c^{2}$ and $\vA^{2} < c^{2}$---both of which hold for
causal equations of state.  Hence, all physical signal speeds in the
relativistic magnetized cylinder are subluminal, guaranteeing a
well-posed initial-value problem for the stability analysis.

%% ================================================================
\section{Astrophysical context: magnetized relativistic jet fragmentation}
\label{sec:rel-12-110-jets}

The results of this section have direct application to the stability
of extragalactic relativistic jets.

\subsection{Fiducial parameters}
\label{sec:rel-12-110-fiducial}

Consider a relativistic jet modelled as a cylinder of (proper-frame)
radius $R_{\mathrm{jet}} \sim 1\,\mathrm{kpc}$, proper energy
density $\varepsilon \sim 10^{-9}\,\mathrm{erg/cm^{3}}$ (dominated by
kinetic energy of relativistic particles), and pressure
$p \sim \varepsilon/3$ (ultra-relativistic equation of state).  A
helical magnetic field threads the jet, with the axial component
$H_{z} \sim 10\,\mu\mathrm{G}$.

For these parameters:
\begin{itemize}
\item The relativistic enthalpy density is
  $w_{\mathrm{tot}} = (\varepsilon + p + H^{2}/4\pi)/c^{2}
  \approx (4\varepsilon/3)/c^{2}$, since $H^{2}/(8\pi) \ll \varepsilon$.
\item The effective gravitational mass density is
  $\rho_{G} = (\varepsilon + 3p)/c^{2} = 2\varepsilon/c^{2}$,
  a factor of 2 larger than the Newtonian-equivalent
  $\varepsilon/c^{2}$.
\item The relativistic Alfv\'en speed is
  $\vA \approx H/\sqrt{4\pi\cdot 4\varepsilon/(3c^{2})}
  = Hc/\sqrt{16\pi\varepsilon/3}$.
\end{itemize}

\subsection{Fragmentation wavelength and growth time}
\label{sec:rel-12-110-fragment}

The minimum wavelength for gravitational instability is
$\lambda_{\min} = 2\pi R/x_{a}$, where $x_{a}$ is determined by
equation~\eqref{eq:rel-12-110-18} with $\rho_{G}$ replacing
$\rho_{0}$.  The relativistic correction \emph{increases} $\rho_{G}$
relative to the Newtonian mass density, thereby strengthening the
gravitational drive and making the jet \emph{more} susceptible to
fragmentation than a purely Newtonian estimate would suggest.  This
effect competes with the enhanced inertia $w_{\mathrm{tot}}$, and the
net result depends on the equation of state:

\begin{itemize}
\item For cold matter ($p \ll \varepsilon$):  the relativistic
  corrections are small, and the classical results of \S~110 apply
  with minor modifications.
\item For ultra-relativistic matter ($p = \varepsilon/3$):
  $\rho_{G} = 2\varepsilon/c^{2}$ and
  $w_{\mathrm{tot}} = 4\varepsilon/(3c^{2})$, so the ratio
  $\rho_{G}/w_{\mathrm{tot}} = 3/2$.  The gravitational drive is
  enhanced by a factor $\rho_{G}/\rho_{0} = 2$ relative to the
  Newtonian case, leading to shorter fragmentation timescales.
\end{itemize}

The exponential sensitivity of $\sigma_{\max}$ to the field
strength (equation~\eqref{eq:rel-12-110-24}) means that even modest
increases in the axial magnetic field can prolong the fragmentation
time by many orders of magnitude---precisely as Chandrasekhar noted
for the non-relativistic case with the galactic-arm parameters of
\S~110.

\subsection{Implications for jet morphology}
\label{sec:rel-12-110-morphology}

The competition between relativistic gravitational fragmentation and
axial magnetic stabilization offers a natural explanation for the
observed diversity of jet morphologies:
\begin{enumerate}
\item \textbf{Knotty jets} (e.g., M87): $H/H_{0,\mathrm{rel}}$
  is moderate, allowing fragmentation at wavelengths of order a few
  $R_{\mathrm{jet}}$, producing the regularly spaced knots observed
  at radio and optical wavelengths.
\item \textbf{Smooth jets} (e.g., Cygnus~A hotspot-dominated):
  $H/H_{0,\mathrm{rel}} > 1$, exponentially suppressing
  fragmentation over the jet's lifetime, so the jet remains coherent
  over hundreds of kiloparsecs.
\item \textbf{Disrupted jets} (e.g., FR~I sources):  the axial field
  is weak relative to $H_{0,\mathrm{rel}}$, and the jet fragments
  rapidly, decelerating and spreading into a plume.
\end{enumerate}

%% ================================================================
\section{Summary and non-relativistic limit}
\label{sec:rel-12-110-summary}

The principal results of this relativistic extension of \S~110 are:

\begin{enumerate}
\item The \textbf{dispersion relation}~\eqref{eq:rel-12-110-15}
  generalizes the Chandrasekhar characteristic
  equation~(126) of \S~110 by replacing $\rho$ with
  appropriate relativistic quantities: $\rho_{G}$ (gravitational
  source, equation~\eqref{eq:rel-12-110-14}) and $w_{\mathrm{tot}}$
  (inertia, equation~\eqref{eq:rel-12-110-3}).

\item The \textbf{relativistic Alfv\'en frequency}
  $\Omega_{a,\mathrm{rel}}$
  (equation~\eqref{eq:rel-12-110-11}) automatically satisfies
  $\vA < c$, ensuring \textbf{causality} of the stabilization
  mechanism.

\item The qualitative conclusion that \emph{no magnetic field can
  fully stabilize the cylinder} remains valid in the relativistic
  regime, owing to the logarithmic divergence of the gravitational
  potential at long wavelengths.

\item For strong fields ($H > H_{0,\mathrm{rel}}$), the exponential
  suppression of the growth rate
  (equation~\eqref{eq:rel-12-110-24}) persists, with the relativistic
  scale set by $\rho_{G}$ and $H_{0,\mathrm{rel}}$.

\item The results provide a framework for understanding the
  \textbf{fragmentation of magnetized relativistic jets} into knots,
  with the jet morphology determined by $H/H_{0,\mathrm{rel}}$.
\end{enumerate}

\paragraph{Non-relativistic limit.}
In the limit $p \to 0$, $\varepsilon \to \rho_{0}c^{2}$,
$H^{2}/(8\pi) \ll \rho_{0}c^{2}$, all relativistic corrections
vanish: $\rho_{G} \to \rho_{0}$, $w_{\mathrm{tot}} \to \rho_{0}$,
$\Omega_{a,\mathrm{rel}} \to \Omega_{a}$,
$H_{0,\mathrm{rel}} \to H_{0}$, and every equation in this section
reduces to its counterpart in Chandrasekhar's \S~110.
